Hyperspectral reflectance spectroscopy enables non-destructive estimation of plant functional traits, yet current deep learning approaches process spectra as one-dimensional sequences, limiting their capacity to capture long-range inter-band dependencies. We investigated whether transforming 1D hyperspectral signals into 2D image representations improves multi-trait retrieval when coupled with convolutional neural networks. Six transformation methods---direct Reshape, Continuous Wavelet Transform, Multi-Window Spectrogram, Two-Dimensional Correlation Spectroscopy, Gramian Angular Fields, and Markov Transition Field---were systematically compared using EfficientNet-B0 on the GreenHyperSpectra dataset (7,897 labeled spectra, eight traits, 400--2450 nm). \jl{We compared 1D vs 2D CNN models trained from scratch, and models using self-supervised learning using a pretrained masked autoencoder (MAE)}. The simple Reshape transformation achieved the best performance \jl{when tuned from scratch }($R^2 = 0.684 \pm 0.001$), significantly improving all other transforms and the state-of-the-art 1D baseline ($R^2 = 0.587$, $+0.097$). For the self-supervised learning, we tested a pretrained masked autoencoder (MAE-2D) on 139,000 unlabeled spectral images. Linear probing of the frozen encoder ($R^2 = 0.646$) exceeded all MAE-1D methods, including the fine-tuned 1D masked autoencoder ($R^2 = 0.645$), demonstrating that 2D encoders learn high-quality spectral representations \jl{even with the encoder kept frozen and no fine-tuning} \sm{NO ES LO MISMO?}. Our results show that the representational advantage of 2D spectral images, rather than architectural complexity or transfer from natural images, drives the performance gain over 1D approaches.

